% !TEX root = ../algebra_02.tex
\section{Действие группы на множестве. Определение и примеры}

\subsection*{Определение}

\begin{Definition}{Действие группы на множество}{}
	Пусть $G$ — группа, $M$ — множество.
	\textbf{Левым действием} группы $G$ на множество $M$ (обозначается $G \curvearrowright M$) называется отображение $G \times M \to M$, заданное правилом $(g, m) \mapsto g \cdot m$ (или просто $gm$), удовлетворяющее двум свойствам:
	\begin{enumerate}
		\item $\forall g_1, g_2 \in G, \forall m \in M \colon (g_1 g_2)m = g_1(g_2 m)$.
		\item $\forall m \in M \colon em = m$, где $e$ — нейтральный элемент группы $G$.
	\end{enumerate}
\end{Definition}

\begin{Remark}{Правое действие}{}
	Аналогично определяется \textbf{правое действие} $M \times G \to M$, где $m \cdot g = mg$, со свойствами: $m(g_1 g_2) = (mg_1)g_2$ и $me = m$.
	Любое правое действие можно превратить в левое, определив $g \cdot m := m g^{-1}$.
\end{Remark}

\subsection*{Примеры действий групп}

\begin{Example}{Примеры}{}
	\begin{enumerate}
		\item Пусть $G = \mathbb{R}$, $M = \mathbb{C}$.
		Действие задается поворотом на угол:
		\[ x \cdot z = z(\cos x + i \sin x) \]
		Проверим свойства:
		\begin{itemize}
			\item $x_1 \cdot (x_2 \cdot z) = z(\cos x_2 + i \sin x_2)(\cos x_1 + i \sin x_1) = z(\cos(x_1+x_2) + i \sin(x_1+x_2)) = (x_1+x_2) \cdot z$.
			\item $0 \cdot z = z(\cos 0 + i \sin 0) = z$.
		\end{itemize}

		\item $G = GL_n(K)$, $M = K^n$.
		Действие — умножение матрицы на вектор.

		\item $G \le S(M)$ (подгруппа симметрической группы).
		Действие задается естественным образом: $g \cdot m = g(m)$.

		\item $G = S_n$, $M = \{1, \dots, n\}^2$.
		Действие: $g \cdot (i, j) = (g(i), g(j))$.

		\item Действие $G$ на себя левыми сдвигами: $M = G$.
		Действие $g \cdot m = gm$.

		\item Действие $G$ на себя сопряжениями: $M = G$.
		Действие $g \cdot m = gmg^{-1}$.
		Проверка свойств:
		\begin{itemize}
			\item $(g_1 g_2) \cdot m = (g_1 g_2)m(g_1 g_2)^{-1} = g_1(g_2 m g_2^{-1})g_1^{-1} = g_1 \cdot (g_2 \cdot m)$.
			\item $e \cdot m = e m e^{-1} = m$.
		\end{itemize}
	\end{enumerate}
\end{Example}
